\documentclass[border=2mm]{standalone}
\usepackage{tikz}

\begin{document}

% Рекурсивная команда для рисования ковра Серпинского
% #1 = x координата левого нижнего угла
% #2 = y координата левого нижнего угла
% #3 = размер квадрата
% #4 = уровень рекурсии (глубина)
\newcommand{\carpet}[4]{%
  \pgfmathtruncatemacro{\level}{#4}%
  \ifnum\level=0
    % Базовый случай: рисуем заполненный квадрат
    \fill[black] (#1,#2) rectangle +(#3,#3);
  \else
    % Делим на 9 частей (3×3)
    \pgfmathsetmacro{\newsize}{#3/3}%
    \pgfmathtruncatemacro{\nextlevel}{\level-1}%
    % Рисуем 8 квадратов (пропускаем центральный)
    \foreach \i in {0,1,2} {
      \foreach \j in {0,1,2} {
        \pgfmathtruncatemacro{\skip}{(\i==1 && \j==1) ? 1 : 0}%
        \ifnum\skip=0
          \pgfmathsetmacro{\newx}{#1 + \i*\newsize}%
          \pgfmathsetmacro{\newy}{#2 + \j*\newsize}%
          \carpet{\newx}{\newy}{\newsize}{\nextlevel}%
        \fi
      }
    }
  \fi
}

\begin{tikzpicture}
  % Рисуем 5 этапов в ряд
  \foreach \stage in {0,1,2,3,4} {
    \begin{scope}[xshift=\stage*4cm]
      % Рамка для наглядности
      \draw[gray, thin] (0,0) rectangle (3,3);
      % Рисуем ковёр
      \carpet{0}{0}{3}{\stage}
      % Подпись
      \node[below] at (1.5,-0.3) {Этап \stage};
    \end{scope}
  }
\end{tikzpicture}

\end{document}
